
\documentclass[11pt,letterpaper]{article}
\input{headings}
\newcommand \recipeName {Pasteis}
\chead{\recipeName}

\begin{document}
\input{title}

Past\'{e}is are typical Brazilian filled deep fried dough pockets. They can be made in various size. When I grew up they were either small like a two-byte snack or a bit bigger like a 3-inch diameter circle folded in half. Currently some of the most popular past\'{e}is in Brazil are found in the popular street farmer markets called "feiras" and thus, are known as "pastel de feira." They are substantially larger and intended as a mid-day snack. There are many different fillings for past\'{e}is, but the most traditional one is ground beef, sometimes augmented with hard-boiled eggs. Cheese filling is also quite popular. For some reason beef and cheese has not caught on. 

I have been frustrated for a long time for not being able to make good dough for past\'{e}is in Canada. My attempts led to a product that was too tough or that did not have enough structure. 

After some research I discovered that the all-purpose flour available in Canada has a higher protein content than the flour available in Brazil. The solution is to mix 20\% of cake flour with 80\% of all-purpose flour. 

The hydration ratio is also a factor. I learned that my early attempts were too wet. The hydration should be very low. At the beginning the dough barely comes together.

Finally,  how much do you knead the dough and how long you let the dough rest play a key role in the final product. In my early attempts, kneading a lot led to tough dough after frying, kneading less lead to little structure to the dough and misshaped past\'{e}is that were hard to handle and easily pierced leading to a mess in the frying pan. I discovered the final process by accident. I had made an attempt that did not work well. In frustration, I left the kitchen and left the dough, covered with plastic, sit on top of the counter at room temperature overnight. The next morning the dough was perfect! Later, in conversation with my Mom, she said that some of the best doughs are the ones that rest overnight at room temperature. Curiously, most of the youtube videos that say that they make a perfect pastel recommend only a 20 minute rest. 

The use of a hard liquor, such as cachaca or vodka is required. All the alcohol evaporates while frying, but it leaves the attractive bubbles in the dough.


I am providing three different fillings. Two traditional ones: beef and cheese. The third is a "novidade" to cater for guests that may be milk intolerant, vegetarian, or both.  In Brazil, people like to follow traditions and also to explore with new things. I developed a spicy vegetable filling.


\vspace{0.3in}

\begin{description}

\item[Dough Ingredients:]\ \\
	\begin{itemize}
	\item 100 grams of cake flour
	\item 400 grams of all-purpose flour
	\item 2 Tablespoon of cooking oil (26 grams)
	\item 200 ml of warm water
	\item 2 Tablespoons of cacha\c{c}a (vodka, or other high proof mild flavour liquor)
	\item 1/2 Tablespoon of salt
	\end{itemize}

\item[Filling Options:]\ \\
	\begin{itemize}
	\item Beef Filling
	\item  \href{ChickenFilling.html}{Chicken Filling}
	\item Cheese filling
	\item Spicy Vegan filling
	\end{itemize}
	
\item[Procedure:]\ \\
	\begin{enumerate}
	\item {\bf Make the dough (one day ahead)}
	\begin{itemize}
	\item Sift both flours into a bowl and mix with a fork.
	\item Drizzle half of the cooking oil over the dough, stir with a fork, drizzle the other half and stir again. The goal is to distribute the oil in the flour without creating a single lump of oil.
	\item In a measuring cup, mix the warm water, cacha\c{c}a and salt, stir well to dissolve the salt.
	\item Add the water mixture to the flour-oil mixture in thirds, stirring in between additions.
	\item Start kneading the dough in the bowl with your hands. It will appear to be quite dry.
	\item Dump the dough into the countertop and keep grabbing the pieces and kneading until it forms a rough ball, it may take several minutes because the dough is quite dry.
	\item Once the dough comes together into a rough lumpy ball, place it back into the bowl, cover with plastic, and let it rest at room temperature for at least one hour.
	\item Move the dough back to the countertop. Spread it out with the heal of your hands into a rectangle and fold so that the seam is at the centre of the rectangle. Repeat: spread into a rectangle and fold, now in the other direction. Do this several times until the dough has a smooth outside texture. The process of stretching and folding will extend the gluten mesh, thus given structure to the dough later.
	\item Place the dough back into the bowl ensuring that it is completely covered by plastic and has no contact with air. 
	\item Let the dough rest overnight at room temperature. At least 12 hours, but it will be good even after 48 hours.
	\end{itemize}
	\item {\bf Roll out the Dough}
	\begin{itemize}
	\item Take a piece of the dough, leaving the remaining dough covered with plastic.
	\item Spread a small amount of flour on the counter and lightly flour the piece of dough on both sides.
	\item If using the rolling pin, roll out the dough to a thin thickness - about 1.0 mm. 
	\item If using the pasta machine, first use a rolling pin to make the dough piece long and thin enough to fit in the machine at the thickiest setting.
	\item Pass the rolled dough through the machine.
	\item Change to the next thinner setting and pass the dough again.
	\item Repeat, stopping at the second-to-last thin setting.
	\item Lay the rolled out piece of dough on a counter top. You should have a long strip of dough that is between 3 and 4 inches wide.
	\end{itemize}
	\item {\bf Make the pasteis}
	\begin{itemize}
	\item Cut the dough strip into squares.
	\item Put a small amount of filling off-center on each square.
	\item Fold the square (you can make rectangles or triangles).
	\item As you fold make sure that you do not trap air into the pastel.
	\item Press all around with your finger to seal.
	\item Using a fork, press all the edges to seal completely.
	\item As you finish, place each pasterl on a baking sheet covered with parchment paper.
	\end{itemize}
	\item {\bf Fry the pasteis}
	\begin{itemize}
	\item Heat neutral cooking oil (I use canola oil) to 360 F.
	\item Place few a pasteis in the hot oil. do not overcrowd the pan or frier.
	\item Cook until lightly brown on one side and then turn over to fry on the other side.
	\item Remove to a wire rack and let cool for a few minutes.
	\item Serve still warm.
	\end{itemize}
	\end{enumerate}
\end{description}
\end{document}



